%<*driver>
\def\nameofplainTeX{plain}
\ifx\fmtname\nameofplainTeX
\input l3docstrip.tex
\preamble

========================================================================
 ltxtalk-theme - themes for ltxtalk

 (C) 2026 Matthias Werner
 E-mail: matthias.werner@informatik.tu-chemnitz.de

 Released under the LaTeX Project Public License v1.3c or later
 See http://www.latex-project.org/lppl.txt
======================================================================== 

\endpreamble
\postamble

 This work is "maintained" (as per LPPL maintenance status) by
 Matthias Werner.
\endpostamble

\generate{
  \file{ltxtalk-theme-tuc-2019.sty}{\from{lttheme-tuc-2019.dtx}{pkg}}
}
  \expandafter\endbatchfile
\fi
 \documentclass[
load-preamble+
]{osgdoc}

\usepackage{enumitem}
\usepackage{microtype}
\usepackage[
languages={de,en},
map={de=german,en=english},
targetlang={job=2},
load babel
]{langselect}

\def\pkgname{ltxtalk-theme-tuc-2019}
\def\thepkg{\pkg{\pkgname}}

\makeindex
\setcnltx{
  package  =  ltxtalk-theme-tuc-2019,
  version  = \UseVersionOf{ltxtalk-theme-tuc-2019.sty},
  date     = \UseDateOf{ltxtalk-theme-tuc-2019.sty},
  title    = {TUC 2019 Theme \ldeen{für}{for} Ltx-Talk},
  info     = {Ltx-Talk\ldeen{ Thema mit dem Corporate
          Design der TU Chemnitz}{ theme with the corporate
          design of the Chemnitz University of Technology}},
  authors  = {Matthias Werner[matthias.werner@informatik.tu-chemnitz.de]},
  abstract={
  \ldeen{
    Dies ist ein Thema für die Präsentationsklasse \cls{ltx-talk} 
    und bildet (in etwa) das Corporate Design 2019 der TU Chemnitz nach. 
    Das Paket benötigt die Themenengine \pkg{ltxtalk-theme}.
    }{
    This is a theme for the \cls{ltx-talk} presentation class 
    and (roughly) replicates the 2019 corporate design of 
    Chemnitz University of Technology. 
    The package requires the \pkg{ltxtalk-theme} theme engine.
    }
  },
  url      =https://github.com/werner-matthias/osglecture,
  build-title
}
\begin{document}

\section{\ldeen{Einführung}{Introduction}}

\ldeen{
Das Theme \code{tuc-2019} bildet das seit 2019 verwendete Corporate
Design der TU Chemnitz für Präsentationen nach. 
}{
% TODO: English translation
}

\begin{sourcecode}
  \usepackage{ltxtalk-theme-tuc-2019}
  \useltxtalktheme{tuc-2019}
\end{sourcecode}

\subsubsection{Aufbau der Folien}

\ldeen{
Auf normalen Folien enthält das grüne Kopfband links das weiße
Universitätslogo. Eine weiße Linie trennt es vom Veranstaltungskontext.
Die erste Textzeile besteht aus Präsentationstitel und aktuellem Abschnitt,
die zweite aus dem aktuellen Unterabschnitt. Die Angaben werden aus
\cs*{title}, \cs*{section} und \cs*{subsection} übernommen.
}{
% TODO: English translation
}

\ldeen{
Der Folientitel und ein optionaler Folienuntertitel stehen im Hauptbereich.
Die Fußzeile ist durch eine dünne Linie abgesetzt und enthält den Autor links,
die Foliennummer in der Seitenmitte und die Webadresse rechts. Alle drei
Felder besitzen den im Corporate Design vorgesehenen Abstand zum Seitenrand.
}{
% TODO: English translation
}

\ldeen{
Die Titelfolie unterscheidet sich bewusst von normalen Folien. Ihr Kopfband
ist höher und zeigt neben dem größeren Universitätslogo den Titel und das
Datum. Im Hauptbereich steht links ein zweites, veranstaltungsbezogenes Logo;
dessen Seitenkanten fluchten mit denen des Universitätslogos. Der Titelblock
beginnt unter dem Text des Kopfbandes. Eine Fußzeile wird auf der Titelfolie
nicht ausgegeben.
}{
% TODO: English translation
}

\subsubsection{Logos und Webadresse}

\ldeen{
Das Universitätslogo, das Logo der Titelfolie und die Webadresse können vor
\cs*{begin}\Marg{document} gesetzt werden. Die Pfade werden wie bei
\cs*{includegraphics} relativ zur Hauptdatei der Präsentation aufgelöst.
Für die mit dem Bundle ausgelieferten Dateien kann die Konfiguration zum
Beispiel so aussehen:
}{
% TODO: English translation
}

\begin{sourcecode}
  \settucthemelogo{images/logos/tuc_white.pdf}
  \settucthemetitlelogo{images/logos/osg.png}
  \settucthemeurl{https://www.tu-chemnitz.de/}
\end{sourcecode}


\subsubsection{Titelfolie}

\ldeen{
Die speziell gestaltete Titelfolie wird mit \cs*{maketuctitle} erzeugt. Der
Befehl legt den benötigten Frame selbst an und darf daher nicht innerhalb
einer \code{frame}-Umgebung stehen. Er verwendet die Angaben aus
\cs*{title}, \cs*{subtitle}, \cs*{author}, \cs*{institute} und \cs*{date}.
Nach der Titelfolie stellt er automatisch die Geometrie normaler Folien
wieder her.
}{
% TODO: English translation
}

\subsubsection{Vollständiges Beispiel}

\begin{example}[code-only]
  \DocumentMetadata{}
  \documentclass{ltx-talk}

  \usepackage{ltxtalk-theme-tuc-2019}
  \settucthemelogo{tuc_white.pdf}
  \settucthemetitlelogo{images/logos/osg.png}
  \settucthemeurl{https://www.tu-chemnitz.de/}
  \useltxtalktheme{tuc-2019}

  \title{Echtzeitsysteme}
  \subtitle{0. Kapitel\\Organisatorisches}
  \author{Prof. Matthias Werner}
  \institute{Professur Betriebssysteme}
  \date{Sommersemester 2026}

  \begin{document}
  \maketuctitle

  \section{Organisatorisches}
  \subsection{Willkommen}
  \begin{frame}
    \frametitle{Willkommen zur Lehrveranstaltung}
    \begin{itemize}
      \item Informationen und Materialien
      \item Termine und Ansprechpartner
    \end{itemize}
  \end{frame}
  \end{document}
\end{example}

\ldeen{
Eine ausführlichere Präsentation befindet sich in
\code{example/example-tuc-2019.tex}.
}{
% TODO: English translation
}



\section{}

\begin{commands}
  \command{settucthemelogo}[\marg{Datei}]
  \ldeen{
    Setzt die Datei für das weiße TU-Chemnitz-Logo im Kopfband. Auf der
    Titelfolie wird es größer dargestellt als auf normalen Folien. Ist die
    Datei nicht auffindbar, gibt das Theme ersatzweise den Text
    \code{TU CHEMNITZ} aus.
  }{
    % TODO: English translation
  }

  \command{settucthemetitlelogo}[\marg{Datei}]
  \ldeen{
    Setzt das veranstaltungsbezogene Logo im linken Bereich der Titelfolie.
    Es wird auf die Breite des TU-Logos skaliert und oben bündig mit dem
    Titelblock angeordnet. Ist die Datei nicht auffindbar, bleibt die
    Logo-Spalte leer.
  }{
    % TODO: English translation
  }

  \command{settucthemetitlelogoraise}[\marg{Länge}]
  \ldeen{
    Verschiebt das sichtbare Logo der Titelfolie um die angegebene Länge nach
    oben, um einen transparenten oberen Bildrand auszugleichen. Der Vorgabewert
    \code{0.21cm} passt zum mitgelieferten OSG-Logo; für eng beschnittene
    Logos ist gewöhnlich \code{0pt} richtig.
  }{
    % TODO: English translation
  }

  \command{settucthemeurl}[\marg{Webadresse}]
  \ldeen{
    Setzt den rechtsbündigen Text in der Fußzeile normaler Folien. Das
    Argument wird unverändert gesetzt; soll die Adresse anklickbar sein, kann
    daher beispielsweise \cs*{url} im Argument verwendet werden.
  }{
    % TODO: English translation
  }

  \command{maketuctitle}
  \ldeen{
    Erzeugt die vollständige Titelfolie des Themes. Der Befehl ist nach
    \cs*{begin}\Marg{document}, aber außerhalb einer \code{frame}-Umgebung
    aufzurufen. Die Titelfolie besitzt keine Fußzeile.
  }{
    % TODO: English translation
  }
\end{commands}


% --------------------------------------------------------------------------
% Implementation
% --------------------------------------------------------------------------

\section{Implementation}

% \AutoDocInput{lttheeme-tuc-2019.dtx}{pkg} % comment out to during text development to speed up latex

\end{document}

%</driver>
%<*pkg>
\NeedsTeXFormat{LaTeX2e}
\ProvidesExplPackage{ltxtalk-theme-tuc-2019}
{2026/07/21}
{0.1}
{TU Chemnitz presentation design introduced in 2019}

\color_set:nnn {tuccolor@black@rgb}{RGB}{0,0,0}
\color_set:nnn{tuccolor@tuc@rgb}{RGB}{0,95,80}
\color_set:nnn{tuccolor@natwi@rgb}{RGB}{112,112,111}
\color_set:nnn{tuccolor@ma@rgb}{RGB}{161,11,112}
\color_set:nnn{tuccolor@mb@rgb}{RGB}{18,51,117}
\color_set:nnn{tuccolor@etit@rgb}{RGB}{228,3,45}
\color_set:nnn{tuccolor@if@rgb}{RGB}{74,130,70}
\color_set:nnn{tuccolor@wiwi@rgb}{RGB}{157,7,54}
\color_set:nnn{tuccolor@phil@rgb}{RGB}{198,83,6}
\color_set:nnn{tuccolor@hsw@rgb}{RGB}{0,117,191}
\color_set:nnn{tuccolor@gold@rgb}{RGB}{141,104,49}
\color_set:nnn{tuccolor@silver@rgb}{RGB}{157,156,156}
\color_set:nnn{tuccolor@info@rgb}{RGB}{85,85,85}
\color_set:nnn{tuccolor@warning@rgb}{RGB}{255,153,0}
\color_set:nnn{tuccolor@danger@rgb}{RGB}{255,0,0}

\color_set:nnn{tuccolor@black@cmyk}{cmyk}{0.00,0.00,0.00,1.00}
\color_set:nnn{tuccolor@tuc@cmyk}{cmyk}{1.00,0.34,0.69,0.30}
\color_set:nnn{tuccolor@natwi@cmyk}{cmyk}{0.00,0.00,0.00,0.70}
\color_set:nnn{tuccolor@ma@cmyk}{cmyk}{0.30,1.00,0.00,0.17}
\color_set:nnn{tuccolor@mb@cmyk}{cmyk}{1.00,0.80,0.00,0.29}
\color_set:nnn{tuccolor@etit@cmyk}{cmyk}{0.00,1.00,0.81,0.00}
\color_set:nnn{tuccolor@if@cmyk}{cmyk}{0.74,0.28,0.85,0.13}
\color_set:nnn{tuccolor@wiwi@cmyk}{cmyk}{0.00,1.00,0.50,0.41}
\color_set:nnn{tuccolor@phil@cmyk}{cmyk}{0.17,0.75,1.00,0.06}
\color_set:nnn{tuccolor@hsw@cmyk}{cmyk}{1.00,0.40,0.00,0.00}
\color_set:nnn{tuccolor@gold@cmyk}{cmyk}{0.29,0.48,0.81,0.36}
\color_set:nnn{tuccolor@silver@cmyk}{cmyk}{0.00,0.00,0.00,0.50}
\color_set:nnn{tuccolor@info@cmyk}{cmyk}{0.00,0.00,0.00,0.76}
\color_set:nnn{tuccolor@warning@cmyk}{cmyk}{0,0.40,1.00,0.00}
\color_set:nnn{tuccolor@danger@cmyk}{cmyk}{0.0,1.00,1.00,0.00}

\tl_new:N \g_ltxtalk_tuc_fak_color
\tl_new:N \g_ltxtalk_tuc_color_space

\msg_new:nnnn { tuc-2019 } { unknown-choice }
  { Unknown~value~'#3'~for~key~'#1'. }
  { Valid~values~are~#2. }

\DeclareKeys{
  fakcolor .choices:nn = {
    tuc, natwi, ma, mb, etit, if, wiwi, phil, hsw
  }{
    \tl_set:Ne \g_ltxtalk_tuc_fak_color \l_keys_choice_str
  },
  fakcolor / unknown .code:n = 
   \msg_error:nneee { tuc-2019 } { unknown-choice }
   { fakcolor } { tuc, natwi, ma, mb, etit, if, wiwi, phil, hsw } % Valid choices
   { \exp_not:n {#1} } % Name of choice key
  ,
  fakcolor .initial:n = tuc,
  colorspace .choices:nn = {
    rgb,cmyk
  }{
    \tl_set:Ne \g_ltxtalk_tuc_color_space \l_keys_choice_str
  },
  colorspace / unknown .code:n =
   \msg_error:nneee { tuc-2019 } { unknown-choice }
   { colorspace } { rgb,cmyk } % Valid choices
   { \exp_not:n {#1} } % Name of choice key
  ,
  colorspace .initial:n = rgb
}
\ProcessKeyOptions\relax
\tl_new:N \g_ltxtalk_tuc_main_color
\exp_args:Nne \tl_set:Nn \g_ltxtalk_tuc_main_color { tuccolor@ \g_ltxtalk_tuc_fak_color @ \g_ltxtalk_tuc_color_space }

\RequirePackage{ltxtalk-theme}

\tl_new:N \g__ltxtalk_tuc_logo_tl
\tl_new:N \g__ltxtalk_tuc_title_logo_tl
\tl_new:N \g__ltxtalk_tuc_url_tl
\dim_new:N \g__ltxtalk_tuc_title_logo_raise_dim
\bool_new:N \g__ltxtalk_tuc_title_bool
\tl_gset:Nn \g__ltxtalk_tuc_logo_tl {tuc_white.pdf}
\tl_gset:Nn \g__ltxtalk_tuc_title_logo_tl {osg.png}
\tl_gset:Nn \g__ltxtalk_tuc_url_tl {https://www.tu-chemnitz.de/}
\dim_gset:Nn \g__ltxtalk_tuc_title_logo_raise_dim {0.21cm}

\NewDocumentCommand \settucthemelogo {m}
  {\tl_gset:Nn \g__ltxtalk_tuc_logo_tl {#1}}
\NewDocumentCommand \settucthemetitlelogo {m}
  {\tl_gset:Nn \g__ltxtalk_tuc_title_logo_tl {#1}}
\NewDocumentCommand \settucthemetitlelogoraise {m}
  {\dim_gset:Nn \g__ltxtalk_tuc_title_logo_raise_dim {#1}}
\NewDocumentCommand \settucthemeurl {m}
  {\tl_gset:Nn \g__ltxtalk_tuc_url_tl {#1}}

\cs_new_protected:Npn \__ltxtalk_tuc_logo:
  {
    \file_if_exist:VTF \g__ltxtalk_tuc_logo_tl
      {
        \includegraphics
          [height=\bool_if:NTF \g__ltxtalk_tuc_title_bool {1.25cm}{0.82cm}]
          {\g__ltxtalk_tuc_logo_tl}
      }
      {\sffamily\bfseries TU~CHEMNITZ}
  }
\cs_new_protected:Npn \__ltxtalk_tuc_title_logo:
  {
    \file_if_exist:VTF \g__ltxtalk_tuc_title_logo_tl
      {
        \includegraphics[width=2.485cm]{\g__ltxtalk_tuc_title_logo_tl}
      }
      {}
  }
\cs_new_protected:Npn \__ltxtalk_tuc_context:
  {
    \sffamily
    \bool_if:NTF \g__ltxtalk_tuc_title_bool
      {\large \@title \\[0.18cm] \normalsize \@date}
      {
        \@title
        \tl_if_blank:VF \l__talk_section_tl
          {~--~\tl_use:N \l__talk_section_tl}
        \tl_if_blank:VF \l__talk_subsection_tl
          {\\\tl_use:N \l__talk_subsection_tl}
      }
  }
\cs_new_protected:Npn \__ltxtalk_tuc_header_line:
  {
    \hbox_to_wd:nn {\paperwidth}
      {
        \skip_horizontal:n {0.4cm}
        \begin{minipage}[t]
          {\bool_if:NTF \g__ltxtalk_tuc_title_bool {2.485cm}{1.63cm}}
          \vspace{0pt}\__ltxtalk_tuc_logo:
        \end{minipage}
        \skip_horizontal:n {0.4cm}
        \begin{minipage}[t]{0.4pt}
          \vspace{0pt}\color{white}
          \rule
            {0.4pt}
            {\bool_if:NTF \g__ltxtalk_tuc_title_bool {1.35cm}{0.82cm}}
        \end{minipage}
        \skip_horizontal:n {0.4cm}
        \begin{minipage}[t]{0.62\paperwidth}
          \vspace{0.08cm}\color{white}\__ltxtalk_tuc_context:
        \end{minipage}
        \hfil
      }
  }
\cs_new_eq:NN \__ltxtalk_tuc_header: \__ltxtalk_tuc_header_line:

\cs_new_protected:Npn \__ltxtalk_tuc_theme_setup:
  {
    \ltxtalksetup
      {
        header = true,
        header-height = 1.45cm,
        left-margin = false,
        right-margin = false,
        footer = true,
        footer-height = 0.52cm
      }
    
    \setltxtalkcolors
      {
        secondary = {HTML}{333333},
        accent = {HTML}{F07D00},
        text = {HTML}{202020},
        background = {HTML}{FFFFFF}
      }

    \exp_args:NnV \color_set:nn {ltxtalk-primary} {\g_ltxtalk_tuc_main_color}  
    \exp_args:NnV  \color_set:nn { structure }   { \g_ltxtalk_tuc_main_color }
    
    \setltxtalktitleplacement{frametitle=content,section=none}

    \defineltxtalkslot{tuc-header}{header}{custom}{role=artifact}
    \setltxtalkslotcontent{tuc-header}{\__ltxtalk_tuc_header:}
    \positionltxtalkslot{tuc-header}{left}{0.38cm}{\paperwidth}{1.05cm}
    \styleltxtalkslot{tuc-header}
      {align=left,color=white,font=\footnotesize\sffamily}

    \defineltxtalkslot{tuc-frametitle}{content}{frametitle}
      {role=heading,heading-level=2}
    \positionltxtalkslot{tuc-frametitle}{left}{top}{\textwidth}{0.9cm}
    \styleltxtalkslot{tuc-frametitle}
      {align=left,color=ltxtalk-primary,font=\Large\sffamily}

    \defineltxtalkslot{tuc-framesubtitle}{content}{framesubtitle}
      {role=heading,heading-level=3}
    \positionltxtalkslot{tuc-framesubtitle}{left}{top}{\textwidth}{0.55cm}
    \styleltxtalkslot{tuc-framesubtitle}
      {align=left,color=ltxtalk-secondary,font=\normalsize\sffamily}

    \defineltxtalkslot{tuc-author}{footer}{custom}{role=note}
    \setltxtalkslotcontent{tuc-author}{\@author}
    \positionltxtalkslot{tuc-author}{0.4cm}{center}{0.36\paperwidth}{0.32cm}
    \styleltxtalkslot{tuc-author}
      {align=left,color=ltxtalk-secondary,font=\tiny\sffamily}

    \defineltxtalkslot{tuc-page}{footer}{custom}{role=navigation}
    \setltxtalkslotcontent{tuc-page}
      {N~--~\@framenumber\ von~\@totalframes}
    \positionltxtalkslot{tuc-page}{center}{center}{0.22\paperwidth}{0.32cm}
    \styleltxtalkslot{tuc-page}
      {align=center,color=ltxtalk-secondary,font=\tiny\sffamily}

    \defineltxtalkslot{tuc-url}{footer}{custom}{role=note}
    \setltxtalkslotcontent{tuc-url}{\tl_use:N \g__ltxtalk_tuc_url_tl}
    \positionltxtalkslot{tuc-url}
      {\dim_eval:n{\paperwidth-0.4cm-0.36\paperwidth}}
      {center}{0.36\paperwidth}{0.32cm}
    \styleltxtalkslot{tuc-url}
      {align=right,color=ltxtalk-secondary,font=\tiny\sffamily}

    \decorateltxtalkarea{header}
      {background=ltxtalk-primary,frame=none,artifact-type=layout}
    \decorateltxtalkarea{footer}
      {
        background=white,
        frame=top,
        rule-color=ltxtalk-secondary,
        rule-thickness=0.35pt,
        artifact-type=layout
      }
  }

\NewDocumentCommand \maketuctitle {}
  {
    \bool_gset_true:N \g__ltxtalk_tuc_title_bool
    \ltxtalksetup
      {header-height=2.4cm,footer=false}
    \positionltxtalkslot{tuc-header}{left}{0.45cm}{\paperwidth}{1.6cm}
    \__ltxtalk_update_text_area:
    \begin{frame}
      \thispagestyle{lttheme}
      \vspace*{0.7cm}
      \noindent
      \skip_horizontal:n {\dim_eval:n{0.4cm-\Gm@lmargin}}
      \begin{minipage}[t]{2.485cm}
        \vspace{0pt}\par
        \vspace{\dim_eval:n{-\g__ltxtalk_tuc_title_logo_raise_dim}}
        \centering\__ltxtalk_tuc_title_logo:
      \end{minipage}
      \skip_horizontal:n {0.4cm+0.4pt+0.4cm}
      \begin{minipage}[t]{0.62\paperwidth}
          \vspace{0pt}
          \sffamily
          \Large\bfseries \@title\par
          \vspace{0.8cm}
          \color{ltxtalk-primary}\large \@subtitle\par
          \vspace{0.75cm}
          \normalsize\color{ltxtalk-secondary}\@author\par
          \small \@institute
      \end{minipage}
    \end{frame}
    \bool_gset_false:N \g__ltxtalk_tuc_title_bool
    \ltxtalksetup
      {header-height=1.45cm,footer=true,footer-height=0.52cm}
    \positionltxtalkslot{tuc-header}{left}{0.38cm}{\paperwidth}{1.05cm}
    \__ltxtalk_update_text_area:
  }

\RegisterLTXTalkTheme{tuc-2019}{\__ltxtalk_tuc_theme_setup:}
%</pkg>
